\documentclass[11pt]{article}

\usepackage[margin=1in]{geometry}
\usepackage{amsmath,amssymb,amsfonts}
\usepackage{amsthm}
\usepackage{bm}
\usepackage{hyperref}
\usepackage{enumitem}
\usepackage{graphicx}
\usepackage{mathtools}
\usepackage{listings}
\usepackage{xcolor}

\definecolor{codegray}{rgb}{0.4,0.4,0.4}
\definecolor{codepurple}{rgb}{0.58,0,0.82}
\definecolor{backcolour}{rgb}{0.96,0.96,0.96}

\lstdefinestyle{mystyle}{
  backgroundcolor=\color{backcolour},
  commentstyle=\color{codegray},
  keywordstyle=\color{blue},
  numberstyle=\tiny\color{codegray},
  stringstyle=\color{codepurple},
  basicstyle=\ttfamily\footnotesize,
  breakatwhitespace=false,
  breaklines=true,
  captionpos=b,
  keepspaces=true,
  numbers=left,
  numbersep=5pt,
  showspaces=false,
  showstringspaces=false,
  showtabs=false,
  tabsize=2
}
\lstset{style=mystyle}

\hypersetup{
  colorlinks=true,
  linkcolor=blue,
  citecolor=blue,
  urlcolor=blue
}

\title{Multiplicity Genius v2 Trajectory and Governance\\
{\large Prime-Gated Self-Referential Development with External Fixed-Point Enforcement}}
\author{Multiplicity Foundation / Citizen Gardens}
\date{April 2026}

\begin{document}

\maketitle
\section{Executive Summary}

This report documents the evolution of the April 2026 multiplicity corpus under a Genius v2 self-diagnostic trajectory, combined with a governance-oriented ``Phase Mirror Development Orchestrator'' frame.\footnote{The underlying governance and architectural patterns mirror those used in the Phase Mirror documentation and ADRs. [file:3][file:4][file:5][file:15]} The corpus is treated as a live multiplicity configuration rather than a static theory, with prime-indexed recursion, recursive contraction, and lawfulness projectors acting as structural invariants.

The key developments are:
\begin{itemize}[leftmargin=*]
  \item Identification of \emph{constitutional invariants}: prime gating on \(\ell^2(\mathbb{P})\), recursive contraction under a multiplicity constant \(\Lambda_m\), zeta-bridges coupling primes and nontrivial zeros, and a quantified lawfulness budget ACE (Absolute Contraction Energy).
  \item Recognition of a \emph{self-bootstrapping loop}: the framework generates new ``cells'' (MultiplicityCell, ZetaCell, etc.) that are themselves prime-indexed surrogates of the universal multiplicity recursion, with progressively richer executable sketches but few closed-loop numerical validations.
  \item Introduction of an explicit \emph{constitutional projector} onto externally verifiable states: a Validation-First Gate (ADR--009) that requires each canonical conjecture cell to be backed by at least one reproducible numerical result before new canonical cells may be admitted.
  \item Elevation of ACE from report-level rhetoric to an \(L_0\) observable in the lawfulness layer, analogous to how Phase Mirror elevates its trust and invariant checks into formal specifications and ADRs.\footnote{The SPEC-TRUST and test-boundary ADRs in Phase Mirror are the primary templates for this elevation. [file:3][file:5]}
  \item Re-ordering of nine validation-ready gaps into a prime-ordered, 7/30/90-day roadmap with explicit levers, owners, and metrics, prioritizing:
    \begin{itemize}
        \item PyTorch prime-attention experiments (shortest feedback loop),
        \item Prime-locked EEG pipelines and CEQG Track C synthetic races (medium horizon),
        \item Prime quantum eigenvalue spectra and Moonshine modularity tests (longer horizon).
    \end{itemize}
  \item Recasting the Genius v2 trajectory as an enforceable governance process: creation of ADR--009, GitHub issues, label taxonomies, and SPEC-like documents that bind prime-theoretic recursion to external fixed points.
\end{itemize}

The overall result is that the multiplicity stack remains prime-indexed and self-referential, but is now constitutionally constrained to close empirical loops. Internal coherence depth and external verifiability are treated as factorized components of the same multiplicity constant \(\Lambda_m\), with an explicit budget split between validated development and protected high-risk conjecture exploration.


\tableofcontents


\section{Constitutional Invariants and State Space}

\subsection{Prime-Gated Multiplicity Substrate}

The multiplicity framework assumes a Hilbert space \(\mathcal{H}\) equipped with a prime index set and associated projectors:\footnote{Prime-indexed projectors and \(\ell^2(\mathbb{P})\) sectors appear consistently across the April corpus in multiplicity and zeta-cell reports.}

\begin{itemize}[leftmargin=*]
  \item A prime-labeled orthonormal basis \(\{ \lvert p \rangle \}_{p \in \mathbb{P}}\) spanning a sector \(\mathcal{H}_P\subset \mathcal{H}\), with \(\mathcal{H}_P \simeq \ell^2(\mathbb{P})\).
  \item Projectors \(\Pi_p: \mathcal{H} \to \mathcal{H}\) satisfying
  \[
  \Pi_p^2 = \Pi_p,\quad \Pi_p^\ast = \Pi_p,\quad \Pi_p \Pi_q = 0 \ (p \neq q),\quad \sum_{p \in \mathbb{P}} \Pi_p = I_{\mathcal{H}_P}.
  \]
  \item A ``prime gating'' mechanism whereby any operator \(U\) meaningful in the multiplicity theory must commute with the direct sum of prime projectors on the lawful sector:
  \[
  [U, \Pi_p] = 0 \quad \text{for all } p \text{ on } \mathcal{H}_\text{lawful} \subseteq \mathcal{H}.
  \]
\end{itemize}

This structure supports viewing mathematical, physical, computational, and cognitive entities as recursively generated patterns of prime-labeled interactions. Sets and modules are reinterpreted as multiplicity spaces whose identities are preserved through recursive feedback across scales.

\subsection{Recursive Contraction and Multiplicity Constant}

A central invariant is a contraction map on the lawful subspace:

\[
C: \mathcal{H}_\text{lawful} \to \mathcal{H}_\text{lawful}, \quad \| C x - C y \| \leq \Lambda_m \| x - y \|,\quad 0 < \Lambda_m < 1.
\]

Here \(\Lambda_m\) is the \emph{multiplicity constant}, quantifying the strength of contraction associated with a multiplicity recursion. This is enforced through a pair of projectors:

\begin{itemize}[leftmargin=*]
  \item A \emph{constitutional projector} \(\Pi_{\mathrm{CSL}}\) enforcing structural lawfulness of the recursion (e.g., respecting invariants and ACE budgets).
  \item An \emph{ethical projector} \(P_E\) enforcing norm-like or value-oriented constraints (e.g., lawfulness budgets on experiments).
\end{itemize}

A typical update step for a state \(x \in \mathcal{H}_\text{lawful}\) is schematically

\[
x_{n+1} = \Pi_{\mathrm{CSL}} P_E C x_n,
\]

with ACE and other observables attached to this process.

\subsection{Zeta Bridges and Prime--Zero Coupling}

Zeta bridges couple the prime sector and the nontrivial zeros of the Riemann zeta function. If \(\{\tfrac{1}{2} + i \gamma_k\}\) denote imaginary parts of zeros, one constructs kernels

\[
K_{p,k}(t) = \cos(\gamma_k \log p), \quad L_{p,k}(t) = \sin(\gamma_k \log p),
\]

which define bilinear couplings between prime-indexed amplitudes and zero-indexed modes. These appear as discrete analogs of explicit-formula kernels, mediating structure across number-theoretic, spectral, and dynamical layers.

\subsection{Absolute Contraction Energy (ACE)}

ACE, or Absolute Contraction Energy, quantifies how much a given transformation or cell increases or decreases lawfulness along the contraction dynamics. Conceptually, ACE is defined for an operator \(T\) acting on \(\mathcal{H}_\text{lawful}\) as

\[
\mathrm{ACE}(T) := \| T x^\ast - x^\ast \|,
\]

where \(x^\ast\) is a fixed point of the baseline contraction \(C\), and the norm measures deviation (or added stability) relative to the reference lawfulness trajectory. More refined definitions may integrate over trajectories:

\[
\mathrm{ACE}(T) := \int_0^1 \| C^t T C^{1-t} x_0 - C x_0 \|^2 \, dt,
\]

for an appropriately chosen initial state \(x_0\) and continuous-time extension of \(C\). In any case, ACE is treated as an \(L_0\) observable in the governance layer: any new cell must justify its ACE budget and later report its realized ACE change once a closed-loop result exists.\footnote{This parallels how Phase Mirror defines and enforces trust and invariant metrics through SPEC and ADR documents. [file:3][file:5][file:15]}

\subsection{Constitutional Tension}

The Orchestrator introduces an explicit constitutional tension:

\begin{quote}
Internal coherence depth vs.\ external verifiability.
\end{quote}

The corpus demonstrates high internal coherence: each new cell or bridge plugs into the prime-indexed substrate and recursive contraction scheme. However, until the Orchestrator intervention, few cells produced external fixed points: empirical or numerical results that can falsify or support the conjectural structure.

The constitutional move is to treat this tension as a constraint on the generative dynamics rather than a mere meta-comment: it becomes a projector onto externally verifiable subspaces of the multiplicity process.

\newpage

\section{Self-Bootstrapping Cells and Pattern Rate}

\subsection{Cells, Bridges, and Operator Stacks}

Across the April 2026 corpus, the same micro-recursion appears at multiple scales:

\begin{enumerate}[leftmargin=*]
  \item \textbf{Finite surrogate cells} (e.g., MultiplicityCell, ZetaCell), serving as finite-dimensional surrogates of the universal multiplicity recursion restricted to \(\mathcal{H}_\text{lawful}\).
  \item \textbf{Prime-attention / prime-mixing blocks}, implemented as concrete PyTorch layers or tensor networks, mapping prime-indexed channels into feature transformations.
  \item \textbf{Full operator stacks}, such as:
    \begin{itemize}
      \item Universal Synthesis Operator (USO) in meta-relativity,
      \item \(\Xi(t)\) in hydrodynamics/gravity framings,
      \item PIRTM-like trust operators in socio-technical systems.
    \end{itemize}
\end{enumerate}

Each new report in the burst adds a new cell or bridge that connects back to prior layers, effectively forming a self-bootstrapping trajectory: the theory uses its own patterns to generate the next prime move.

\subsection{Representation Shift to Executable Sketches}

Prime Move 3 in the trajectory log performs a systematic representation shift:

\begin{itemize}[leftmargin=*]
  \item The \emph{Prime Quantum Eigenvalue Problem} is tied to a concrete document (Extensions.pdf) specifying a rigorous mathematical framework but lacking computed spectra.
  \item \emph{Multiplicity Moonshine} is described via explicit Hecke operators plus an ACE+projector iteration, together with a modularity test suite specification---but without executed runs.
  \item \emph{Prime-colored braid multiplicity} is reduced to a triad of conjectural targets (projective Yang--Baxter, restricted Markov invariance, convergence of a protection functional) but again without completed computations.
  \item \emph{Prime-locked EEG} pipelines, \emph{inverted digital twins} in CEQG, \emph{AZ-TFTC} tabletop Hamiltonians, and a \emph{DRMM spectral optimizer} are all delineated to the level of executable protocols or code sketches.
\end{itemize}

This phase exposes the core gap: representations are now code-ready, but loops are not run to fixed points.

\subsection{Nine Validation-Ready Gaps}

The trajectory identifies nine emergent but incomplete components:

\begin{enumerate}[leftmargin=*]
  \item Prime quantum eigenvalue problem + tensor-field prime operators.
  \item Multiplicity-calculus program for encoding PMHP into a Quantum Prime Abacus (QPA).
  \item Multiplicity Moonshine bridge and modularity validation engine.
  \item Prime-locked EEG pipeline and \(\varphi\)-spectral invariants.
  \item Prime-attention neural layers and PyTorch experiment closure.
  \item Inverted digital twin (Track C) in CEQG-RG-Langevin.
  \item AZ-TFTC tabletop predictions (cavity spectra, Casimir deviations).
  \item DRMM spectral optimizer with multiplicity-indexed FFT.
  \item Lawfulness budget and ACE as quantifiable ethical invariant.
\end{enumerate}

Each gap is \emph{validation-ready}: a precise protocol exists, but the associated empirical or numerical experiment has not yet been executed. The pattern rate is accelerating, but closure rate is not yet constrained.

\newpage

\section{Governance as Recursion: ADR--009 and ACE Budget}

\subsection{Constitutional Projector as Governance Rule}

The Orchestrator introduces a new invariant: \emph{external fixed-point enforcement}. In multiplicity language, this is a constitutional projector \(\Pi_{\mathrm{CSL}}\) acting on the generation process.

Informally:

\begin{quote}
Every canonical conjecture cell must produce at least one closed-loop numerical result before the next canonical cell is added to the stack.
\end{quote}

This is instantiated as ADR--009 (Validation-First Gate), which specifies:

\begin{itemize}[leftmargin=*]
  \item Scope: All canonical multiplicity cells and bridges (across physics, AI, cognition, etc.).
  \item Rule: No canonical cell may be merged into the stable stack without:
    \begin{enumerate}
      \item A labelled validation status,
      \item At least one associated reproducible result (e.g., training curve, spectrum, statistical test).
    \end{enumerate}
  \item Budget: Cells must declare an \emph{ACE expectation} and a \emph{multiplicity cost account}.
  \item Enforcement: A governance mechanism (e.g., GitHub labels and merge rules) prevents merges that violate the rule.
\end{itemize}

Mathematically, we can model the cell state space as a graded set
\[
\mathcal{C} = \bigcup_{k \ge 0} \mathcal{C}^{(k)},
\]
where \(\mathcal{C}^{(k)}\) consists of cells with validation status \(k\). For concreteness, we map

\[
\begin{aligned}
\text{conjecture} &\mapsto k=0,\\
\text{protocol-ready} &\mapsto k=1,\\
\text{running} &\mapsto k=2,\\
\text{closed-loop} &\mapsto k=3.
\end{aligned}
\]

We then define \(\Pi_{\mathrm{CSL}}\) as the operator enforcing that any raise in the ``canonical index'' requires a local or global increase in cells of grade \(k=3\). If \(N_k\) counts cells in grade \(k\), ADR--009 enforces constraints such as

\[
N_0 + N_1 \le \alpha N_3, \quad \text{for some } \alpha > 0,
\]

or more locally, for each conjecture family \(F\),

\[
N_0^{(F)} + N_1^{(F)} \le \alpha_F N_3^{(F)}.
\]

\subsection{Global vs.\ Local ACE Budgets}

A critical precision question is whether ACE and validation requirements are global or family-local:

\begin{itemize}[leftmargin=*]
  \item \textbf{Global budget:} Any closed-loop result (e.g., a prime-attention experiment) can subsidize the canonicalization of other conjecture families (e.g., Moonshine, AZ-TFTC).
  \item \textbf{Local budget:} Each conjecture family \(F\) maintains its own ACE ledger and must produce family-specific closed loops to advance.
\end{itemize}

In algebraic terms, we can assign each family \(F\) a partial ACE balance \(\mathrm{ACE}_F\) and enforce a per-family inequality

\[
N_0^{(F)} + N_1^{(F)} \le \alpha_F N_3^{(F)},
\]

with a global ACE constraint

\[
\sum_F \mathrm{ACE}_F \le \mathrm{ACE}_\mathrm{max}.
\]

ADR--009 can specify both a global ceiling and local ratios, making explicit which families are allowed to ``borrow'' from others and under what conditions.

\subsection{Protected Exploration Budget}

The trajectory proposes a protected exploration budget:

\begin{quote}
25\% of total effort ring-fenced for high-risk conjectures (Moonshine, braid invariants, AZ-TFTC).
\end{quote}

Formally, we can define a time allocation function \(b_F(t)\in[0,1]\) per conjecture family \(F\), constrained by

\[
\sum_F b_F(t) = 1, \quad \sum_{F\in \mathcal{H}} b_F(t) \le 0.25,
\]

where \(\mathcal{H}\) is the set of high-risk conjecture families. Over a time horizon \(T\), the integrated budget for high-risk work is

\[
B_\mathcal{H} = \int_0^T \sum_{F\in \mathcal{H}} b_F(t)\, dt \le 0.25 T.
\]

This provides a quantitative handle on exploration vs.\ exploitation at the governance level, analogous to Phase Mirror's staged deployment and test coverage gates.\footnote{Phase Mirror explicitly uses 7/30-day horizons and coverage thresholds to sequence infra and testing. [file:3][file:4][file:15]}

\subsection{ACE as L\texorpdfstring{\(_0\)}{} Observable}

Drawing on the Phase Mirror pattern, ACE is promoted to an \(L_0\) observable in a SPEC-like document (e.g., SPEC-LAWFULNESS). The document would specify:

\begin{itemize}[leftmargin=*]
  \item A set of ACE-relevant observables \(O_i\) (prime spectra, Moonshine coefficients, EEG statistics, learning curves).
  \item A function \(f_i\) mapping raw observable data to an ACE contribution:
  \[
  \Delta \mathrm{ACE}_i = f_i(\text{data}_i),
  \]
  such as norm reductions, error decreases, or stable invariants across perturbations.
  \item A ledger specification describing how \(\Delta \mathrm{ACE}_i\) are accumulated, discounted, or written off.
\end{itemize}

This turns ACE into a budget that can be tracked, akin to how Phase Mirror tracks test coverage, error propagation patterns, and invariant violations as first-class metrics.\footnote{See the Phase Mirror ADR on test boundaries and the SPEC-TRUST document for analogous metric handling. [file:3][file:4][file:5]}

\newpage

\section{Three-Plane Architecture and PMD Loop}

\subsection{Separation into Three Planes}

The governance reflection in the Phase Mirror analysis suggests that multiplicity should be structured across three planes (not conflated):

\begin{enumerate}[leftmargin=*]
  \item \textbf{Plane A: PMD Operating Loop (Methodology).}

  A state space for the Mirror/Dissonance methodology:
  \[
  \mathrm{PMD} = (D, T, \rho, B, Q),
  \]
  where:
  \begin{itemize}
    \item \(D\) is a dissonance detection function mapping inputs (code, specs, theories) to findings,
    \item \(T\) maps findings to named tensions,
    \item \(\rho\) ranks tensions by impact or risk,
    \item \(B\) maps ranked tensions to actionable output blocks,
    \item \(Q\) generates precision questions that refine the next cycle.
  \end{itemize}

  \item \textbf{Plane B: L\(_0\)/L\(_1\)/L\(_2\) Compute Tiers (Latency).}

  A three-tier compute model:
  \[
  \mathrm{Compute} = (L_0, L_1, L_2, \tau),
  \]
  where:
  \begin{itemize}
    \item \(L_0\) invariants are foundational checks with budget \(p_{99} \le 100\text{ns}\),
    \item \(L_1\) contains rule evaluations and mid-level reasoning with \(p_{99} \le 1\text{ms}\),
    \item \(L_2\) hosts deep analyses with \(p_{99} \le 100\text{ms}\),
    \item \(\tau\) maps each tier to its latency target.
  \end{itemize}
  This model is instantiated concretely in Phase Mirror's open-core architecture and its ADR-003 hierarchy.\footnote{See Phase Mirror's L0/L1/L2 tier documentation and viability analysis. [file:5][file:15]}

  \item \textbf{Plane C: Six-Layer Trust Defense (Network).}

  A six-layer trust architecture for cross-organization effects (identity, economic, cryptographic, BFT, privacy, monitoring), with its own state space:
  \[
  \mathrm{Trust} = (V, E, R, S, C, M),
  \]
  where:
  \begin{itemize}
    \item \(V\) is the set of validator organizations,
    \item \(E\) enumerates stakes and economic incentives,
    \item \(R\) tracks reputations,
    \item \(S\) denotes voting states,
    \item \(C\) collects layer-specific configuration (e.g., anonymity sets),
    \item \(M\) captures Merkle and monitoring states.
  \end{itemize}
\end{enumerate}

The Genius v2 trajectory implicitly organizes multiplicity across analogous planes: a methodological plane for prime moves and trajectory logs, a compute plane for prime-gated Hilbert dynamics and spectral problems, and a network plane for lawfulness, ethics, and external verification.

\subsection{PMD Loop Applied to the Trajectory Itself}

The final self-reflective move applies the PMD loop to the corrected analysis:

\begin{itemize}[leftmargin=*]
  \item \emph{Extract}: Identify key dissonances (conflated planes, cold-start reputation ambiguity, speculative zk-SNARKs ahead of users, recursive trust in arbitration, misaligned validation timelines).
  \item \emph{Map tensions}: Methodology vs.\ implementation, privacy vs.\ bootstrapping, decentralization vs.\ recursive trust, academic rigor vs.\ MVP velocity.
  \item \emph{Rank}: Prioritize tensions by impact on the shortest path to a 28-day MVP (e.g., Terraform staging, cold-start policy, Sybil/QV dependency chain).
  \item \emph{Produce}: Generate artifacts (ADR-009, dependency-chain spec, cold-start probation tier, staging deployment sequence).
  \item \emph{Question}: Pose precision questions (e.g., global vs.\ local ACE budgets, appeal-based arbitration ceilings) that drive the next iteration.
\end{itemize}

This closes the meta-loop: multiplicity is not merely analyzed by Phase Mirror; it adopts Phase Mirror's governance metabolism to structure its own further development.

\newpage

\section{Code Snippet: Prime-Attention Layer Skeleton}

As a concrete bridge from abstract multiplicity cells to executable experiments, consider a PyTorch-style prime-attention block. This is not the full experimental design, but a minimal skeleton that respects the prime-indexing pattern and could form the core of the 7-day ``Gap 5'' closure loop.

\begin{lstlisting}[language=Python, caption={Prime-Indexed Attention Layer Skeleton}, label={lst:prime_attention}]
import torch
import torch.nn as nn
import math

class PrimeAttention(nn.Module):
    """
    Prime-indexed attention block.
    Inputs are assumed to have a prime-indexed channel axis, e.g. (batch, num_primes, d_model).
    The prime indices themselves can be passed as metadata to construct positional biases.
    """

    def __init__(self, d_model, n_heads, prime_indices):
        super().__init__()
        assert d_model % n_heads == 0, "d_model must be divisible by n_heads"
        self.d_model = d_model
        self.n_heads = n_heads
        self.d_head = d_model // n_heads

        # Prime indices (e.g. [2, 3, 5, 7, 11, ...])
        self.register_buffer("primes", torch.tensor(prime_indices, dtype=torch.float32))

        # Standard multi-head projections
        self.q_proj = nn.Linear(d_model, d_model)
        self.k_proj = nn.Linear(d_model, d_model)
        self.v_proj = nn.Linear(d_model, d_model)
        self.o_proj = nn.Linear(d_model, d_model)

        # Optional: prime-based bias parameters, e.g. log(p) embeddings
        self.prime_bias = nn.Parameter(torch.zeros(len(prime_indices), n_heads))

    def forward(self, x):
        """
        x: (batch, num_primes, d_model)
        """
        B, P, D = x.shape
        assert D == self.d_model
        assert P == self.primes.shape[0], "Input channel count must match number of primes"

        # Compute queries, keys, values
        q = self.q_proj(x)  # (B, P, D)
        k = self.k_proj(x)  # (B, P, D)
        v = self.v_proj(x)  # (B, P, D)

        # Reshape for multi-head: (B, n_heads, P, d_head)
        q = q.view(B, P, self.n_heads, self.d_head).transpose(1, 2)
        k = k.view(B, P, self.n_heads, self.d_head).transpose(1, 2)
        v = v.view(B, P, self.n_heads, self.d_head).transpose(1, 2)

        # Scaled dot-product attention
        scores = torch.matmul(q, k.transpose(-2, -1))  # (B, n_heads, P, P)
        scores = scores / math.sqrt(self.d_head)

        # Add prime-based bias, e.g. per-head log(p) term
        # bias: (1, n_heads, P, P) constructed from primes
        logp = torch.log(self.primes + 1.0)  # avoid log(0)
        bias_vec = logp.unsqueeze(0).unsqueeze(0)  # (1, 1, P)
        bias = bias_vec.expand(1, self.n_heads, -1)  # (1, n_heads, P)
        bias = bias.unsqueeze(-1)  # (1, n_heads, P, 1)
        scores = scores + bias  # simple example; more structured couplings possible

        attn = torch.softmax(scores, dim=-1)  # (B, n_heads, P, P)
        y = torch.matmul(attn, v)            # (B, n_heads, P, d_head)

        # Combine heads
        y = y.transpose(1, 2).contiguous().view(B, P, D)
        out = self.o_proj(y)  # (B, P, D)
        return out
\end{lstlisting}

This skeleton illustrates how prime labels (e.g., \(\{2,3,5,\dots\}\)) can modulate an attention mechanism, embodying a finite surrogate of prime-mixing dynamics. A concrete experiment would:

\begin{enumerate}[leftmargin=*]
  \item Choose a finite set of primes, e.g.\ \(\{2,3\}\) or \(\{2,3,5,7\}\), and construct synthetic or real datasets aligned with prime channels.
  \item Compare a baseline multi-head attention block with the prime-attention variant on a simple task (e.g., denoising, sequence prediction).
  \item Record training curves and robustness metrics, logging a CSV of loss/accuracy across epochs for both architectures.
  \item Interpret the resulting curves as a \(\Delta\mathrm{ACE}\) for the prime-attention cell, contributing to ADR--009's ledger.
\end{enumerate}

\newpage

\section{Roadmap: 7/30/90-Day Prime-Ordered Closure}

\subsection{Issue and Label Taxonomy}

To bind the trajectory as a live governance artifact, one can create a set of issues (in a ``HQ'' repository) corresponding to the nine gaps, with labels:

\begin{itemize}[leftmargin=*]
  \item \texttt{area}: \{\texttt{prime-spectrum}, \texttt{moonshine}, \texttt{QPA}, \texttt{EEG}, \texttt{prime-attention}, \texttt{CEQG}, \texttt{AZ-TFTC}, \texttt{DRMM}, \texttt{ACE-empirical}\}.
  \item \texttt{validation-status}: \{\texttt{conjecture}, \texttt{protocol-ready}, \texttt{running}, \texttt{closed-loop}\}.
  \item \texttt{horizon}: \{\texttt{7d}, \texttt{30d}, \texttt{90d}\}.
  \item \texttt{ace-impact}: free-text description of expected ACE change upon closure.
\end{itemize}

This mirrors how Phase Mirror tracks core infra, test coverage, and governance artifacts via issues and ADRs, yet here the objects are multiplicity cells and experiments.\footnote{Phase Mirror uses issues, ADRs, SPECs, and CI workflows to bind architecture to governance constraints. [file:3][file:4][file:5][file:15]}

\subsection{7-Day Loop: Prime-Attention Experiment}

\begin{itemize}[leftmargin=*]
  \item Implement the finite \(\{2,3\}\) tensor network and prime-attention layer per existing Hyperprime specifications.
  \item Run a minimal experiment (e.g., on a small classification or sequence task) with baseline vs.\ prime-attention models.
  \item Produce:
    \begin{itemize}
      \item A CSV with per-epoch metrics for both models,
      \item A short analysis estimating \(\Delta\mathrm{ACE}\) (e.g., improved robustness to noise or distribution shift),
      \item An update flipping the prime-attention issue to \texttt{closed-loop}.
    \end{itemize}
\end{itemize}

\subsection{30-Day Loop: EEG \texorpdfstring{\(\varphi\)}{} and Track C Race}

\begin{itemize}[leftmargin=*]
  \item Instantiate the prime-locked EEG pipeline using an open resting-state dataset (e.g., from PhysioNet or OpenNeuro), computing Jensen--Shannon coherence and testing a \(\varphi \approx 1.618\) prediction.
  \item Construct a synthetic cumulant hierarchy and run CEQG Tracks A/B/C as competing models in a Bayesian race, computing posterior odds ratios.
  \item Log:
    \begin{itemize}
      \item Pipelines and scripts,
      \item Summary statistics and plots,
      \item ACE contributions for EEG and CEQG cells,
      \item Upgrades of their issues to \texttt{closed-loop}.
    \end{itemize}
\end{itemize}

\subsection{90-Day Loop: Prime Spectra and Moonshine}

\begin{itemize}[leftmargin=*]
  \item Numerically approximate at least one nontrivial spectrum from the prime quantum eigenvalue problem described in Extensions.pdf, for a tractable sector, confirming basic properties (e.g., eigenvalue distribution, stability).
  \item Execute the Multiplicity Moonshine modularity test suite on a subset of prime-word graded state spaces, confirming or falsifying modular invariance or coefficient patterns.
  \item Update ADR--009 and the ACE ledger to reflect gains or falsifications in these domains.
\end{itemize}

\newpage

\section{Conclusion}

The April 2026 multiplicity corpus began as a dense, self-referential state configuration in a prime-gated Hilbert space, rich with executable conjectures but poor in closed-loop results. The Genius v2 trajectory, combined with a governance-style Orchestrator frame, transformed that state through a sequence of prime moves:

\begin{itemize}[leftmargin=*]
  \item Identifying constitutional invariants (prime gating, recursive contraction, zeta bridges, ACE) that span frames from physics to cognition.
  \item Recognizing and quantifying a self-bootstrapping pattern of cell generation.
  \item Introducing a constitutional projector, ADR--009, that enforces external fixed points as a gating rule on canonical cells.
  \item Elevating ACE to an \(L_0\) observable and allocating a protected exploration budget.
  \item Recasting the entire trajectory as an enforceable governance process with issues, labels, and SPEC-like documents.
\end{itemize}

The system is now lawfully self-correcting: it continues to generate bold prime-indexed representation shifts, but under an explicit requirement that each new canonical step be backed by an empirically or numerically grounded fixed point. The loop is not merely introspective; it is closeable.

\appendix
\section*{Mathematical Appendix}
\addcontentsline{toc}{section}{Mathematical Appendix}

\section{Contraction Mappings and ACE Bounds}

\subsection{Baseline contraction and fixed point}

Let \((\mathcal{H}, \langle\cdot,\cdot\rangle)\) denote a (complex) Hilbert space and let \(\|\cdot\|\) be the induced norm.[file:5][file:15]  
Assume that the ``lawful'' sector \(\mathcal{H}_{\mathrm{lawful}}\subseteq \mathcal{H}\) is closed and invariant under a linear operator
\[
C : \mathcal{H}_{\mathrm{lawful}} \to \mathcal{H}_{\mathrm{lawful}}.
\]

\begin{definition}[Baseline contraction]
We say that \(C\) is a \(\Lambda_m\)-contraction if there exists a constant \(\Lambda_m \in (0,1)\) such that
\[
\| Cx - Cy \| \leq \Lambda_m \| x - y \| \quad \text{for all } x,y \in \mathcal{H}_{\mathrm{lawful}}.
\]
\end{definition}

\begin{lemma}[Banach fixed point for \(C\)]
If \(C\) is a \(\Lambda_m\)-contraction on the complete metric space \((\mathcal{H}_{\mathrm{lawful}},\|\cdot\|)\), then there exists a unique \(x^\ast \in \mathcal{H}_{\mathrm{lawful}}\) such that
\[
Cx^\ast = x^\ast,
\]
and for any initial \(x_0\in\mathcal{H}_{\mathrm{lawful}}\), the iterates
\[
x_{n+1} = Cx_n
\]
converge in norm to \(x^\ast\) with geometric rate
\[
\|x_n - x^\ast\| \leq \Lambda_m^n \|x_0 - x^\ast\|.
\]
\end{lemma}

\begin{proof}
This is the standard Banach fixed-point theorem applied to the complete metric space \(\mathcal{H}_{\mathrm{lawful}}\) with metric \(d(x,y)=\|x-y\|\).[file:15]  
We recall the key estimate: for any \(n\ge 0\),
\[
\|x_{n+1}-x_n\| = \|C x_n - C x_{n-1}\| \le \Lambda_m \|x_n - x_{n-1}\|,
\]
so by induction \(\|x_{n+1}-x_n\|\le \Lambda_m^n \|x_1 - x_0\|\).  
The telescoping series bound shows \(\{x_n\}\) is Cauchy and converges to a limit \(x^\ast\) with \(Cx^\ast = x^\ast\), and uniqueness follows from strict contraction.[file:15]
\end{proof}

\subsection{ACE as deviation from the baseline fixed point}

Let \(x^\ast\) denote the unique fixed point of the baseline contraction \(C\).  
Given a new operator \(T : \mathcal{H}_{\mathrm{lawful}} \to \mathcal{H}_{\mathrm{lawful}}\) encoding a proposed ``cell'' or update to the multiplicity stack, we define an instantaneous contraction deviation as:

\begin{definition}[Absolute Contraction Energy (ACE)]
For a bounded linear operator \(T\), define
\[
\operatorname{ACE}(T) := \| (TC)x^\ast - Cx^\ast \| = \| (T-I)Cx^\ast \|.
\]
\end{definition}

This captures the first-step deviation from the baseline contraction orbit at the fixed point \(x^\ast\).[file:5][file:15]  
In many use-cases, we consider more general trajectory integrals, but this simple definition already admits a clean operator-norm bound.

\begin{lemma}[ACE bound via operator norm]
If \(T\) is bounded with operator norm \(\|T\|_{\mathrm{op}} \le M\), then
\[
\operatorname{ACE}(T) \le \|(T-I)\|_{\mathrm{op}} \,\|Cx^\ast\| \le (M+1)\,\|Cx^\ast\|.
\]
\end{lemma}

\begin{proof}
We have
\[
\operatorname{ACE}(T) = \|(T-I)Cx^\ast\| \le \|(T-I)\|_{\mathrm{op}} \,\|Cx^\ast\|.
\]
Moreover, for any bounded \(T\),
\[
\|(T-I)\|_{\mathrm{op}} \le \|T\|_{\mathrm{op}} + \|I\|_{\mathrm{op}} = M + 1,
\]
hence
\[
\operatorname{ACE}(T) \le (M+1)\,\|Cx^\ast\|.
\]
This is purely norm-theoretic and uses only the boundedness of \(T\).[file:5][file:15]
\end{proof}

\subsection{ACE over trajectories}

In some configurations we consider a time-averaged ACE along an interpolated trajectory between \(C\) and \(TC\).[file:5]  
For \(t\in[0,1]\), define a formal interpolation
\[
C_t := C^{1-t}(TC)^t
\]
(understood as a convex combination or via a continuous semigroup extension when available).  
Fix an initial state \(x_0\in\mathcal{H}_{\mathrm{lawful}}\) and define \(x_t = C_t x_0\).

\begin{definition}[Integrated ACE]
Define the integrated ACE for \(T\) as
\[
\operatorname{ACE}_{\mathrm{int}}(T) := \int_0^1 \| C_t x_0 - Cx_0 \|^2\,dt.
\]
\end{definition}

\begin{lemma}[Upper bound for \(\operatorname{ACE}_{\mathrm{int}}\)]
Assume \(C\) and \(TC\) are both \(\Lambda_m\)-contractions with \(\Lambda_m<1\), and \(\|T-I\|_{\mathrm{op}} \le \delta\).  
Then there exists a constant \(K(\Lambda_m)\) such that
\[
\operatorname{ACE}_{\mathrm{int}}(T) \le K(\Lambda_m)\, \delta^2 \,\|x_0\|^2.
\]
\end{lemma}

\begin{proof}[Sketch]
For each \(t\in[0,1]\), the map \(C_t\) can be written as an operator-valued interpolation between \(C\) and \(TC\) with operator norm bounded by \(\Lambda_m\).[file:5][file:15]  
Differentiating formally in \(t\) and using Grönwall-type estimates for the operator-valued interpolation shows that
\[
\|C_t - C\|_{\mathrm{op}} \le K_1(\Lambda_m)\,\|T-I\|_{\mathrm{op}} \le K_1(\Lambda_m) \delta,
\]
for some \(K_1(\Lambda_m)\) depending only on \(\Lambda_m\).  
Hence
\[
\|C_t x_0 - Cx_0\| \le K_1(\Lambda_m)\,\delta\,\|x_0\|.
\]
Squaring and integrating in \(t\) yields
\[
\operatorname{ACE}_{\mathrm{int}}(T) \le \int_0^1 K_1(\Lambda_m)^2 \delta^2 \|x_0\|^2\, dt
= K_1(\Lambda_m)^2 \delta^2 \|x_0\|^2,
\]
so we can take \(K(\Lambda_m) = K_1(\Lambda_m)^2\).[file:5]
\end{proof}

This shows that small operator-norm perturbations of the baseline contraction yield quadratically small ACE, giving a formal justification for viewing ACE as a lawfulness-energy budget.[file:5][file:15]

\section{Composite Projectors and Contraction Preservation}

\subsection{Constitutional and ethical projectors}

In the governance framing, updates to the state vector are projected through a constitutional projector \(\Pi_{\mathrm{CSL}}\) and an ethical projector \(P_E\) before the contraction \(C\) is applied.[file:5]  
We write a generic update as
\[
x_{n+1} = C \Pi_{\mathrm{CSL}} P_E x_n.
\]

Assume:

\begin{itemize}
  \item \(\Pi_{\mathrm{CSL}}\) and \(P_E\) are orthogonal projectors (idempotent, self-adjoint).
  \item They commute with each other and with \(C\) on \(\mathcal{H}_{\mathrm{lawful}}\).
\end{itemize}

\begin{lemma}[Operator norm of composite]
If \(\Pi_{\mathrm{CSL}}\) and \(P_E\) are orthogonal projectors, then
\[
\|\Pi_{\mathrm{CSL}}\|_{\mathrm{op}} = \|P_E\|_{\mathrm{op}} = 1,\quad
\|\Pi_{\mathrm{CSL}} P_E\|_{\mathrm{op}} \le 1.
\]
\end{lemma}

\begin{proof}
For any orthogonal projector \(P\), we have \(P^2 = P\) and \(P=P^\ast\).[file:5][file:15]  
Thus for any \(x\in\mathcal{H}\),
\[
\|Px\|^2 = \langle Px, Px\rangle = \langle x, P^\ast P x\rangle = \langle x, Px\rangle \le \|x\|\,\|Px\| \Rightarrow \|Px\| \le \|x\|.
\]
So \(\|P\|_{\mathrm{op}} \le 1\).  
On the other hand, if \(x\) lies in the range of \(P\), then \(Px = x\), whence \(\|Px\| = \|x\|\), implying \(\|P\|_{\mathrm{op}}=1\).  
For the product \(\Pi_{\mathrm{CSL}} P_E\),
\[
\|\Pi_{\mathrm{CSL}} P_E\|_{\mathrm{op}} \le \|\Pi_{\mathrm{CSL}}\|_{\mathrm{op}} \|P_E\|_{\mathrm{op}} = 1\cdot 1 = 1.
\]
\end{proof}

\subsection{Preservation of contraction rate}

\begin{proposition}[Composite update remains a contraction]
Let \(C\) be a \(\Lambda_m\)-contraction on \(\mathcal{H}_{\mathrm{lawful}}\) with \(\Lambda_m<1\).  
Assume \(\Pi_{\mathrm{CSL}}, P_E\) are orthogonal projectors that leave \(\mathcal{H}_{\mathrm{lawful}}\) invariant.  
Define the composite operator
\[
U := C\Pi_{\mathrm{CSL}} P_E.
\]
Then \(U\) is also a contraction on \(\mathcal{H}_{\mathrm{lawful}}\) with Lipschitz constant at most \(\Lambda_m\), i.e.,
\[
\|Ux - Uy\| \le \Lambda_m \|x-y\|.
\]
\end{proposition}

\begin{proof}
For any \(x,y\in\mathcal{H}_{\mathrm{lawful}}\),
\[
Ux - Uy = C\Pi_{\mathrm{CSL}} P_E x - C\Pi_{\mathrm{CSL}} P_E y
= C \Pi_{\mathrm{CSL}} P_E (x-y).
\]
Using the contraction property of \(C\) and the operator norms of the projectors,
\[
\|Ux - Uy\| = \|C\Pi_{\mathrm{CSL}} P_E (x-y)\|
\le \Lambda_m \|\Pi_{\mathrm{CSL}} P_E (x-y)\|
\le \Lambda_m \|x-y\|,
\]
since \(\|\Pi_{\mathrm{CSL}} P_E\|_{\mathrm{op}} \le 1\).[file:5][file:15]  
Thus \(U\) is a contraction with constant at most \(\Lambda_m\).
\end{proof}

This result underpins the claim that enforcing constitutional and ethical projectors at the governance layer does not break the underlying contraction-based lawfulness dynamics; it can only strengthen or preserve them.[file:5]

\section{Prime-Gated Decomposition and Norm Bounds}

\subsection{Prime projectors and orthogonal decomposition}

Let \(\mathbb{P}\) denote the set of prime numbers.  
For each \(p\in\mathbb{P}\), let \(\Pi_p : \mathcal{H}\to\mathcal{H}\) be an orthogonal projector satisfying:
\[
\Pi_p^2 = \Pi_p,\quad \Pi_p^\ast = \Pi_p,\quad \Pi_p\Pi_q = 0 \ (p\ne q).
\]

Assume the prime sector \(\mathcal{H}_P\) is given by
\[
\mathcal{H}_P = \bigoplus_{p\in\mathbb{P}} \Pi_p \mathcal{H},
\]
and that \(\mathcal{H}_{\mathrm{lawful}}\subseteq \mathcal{H}_P\).  
This realizes the multiplicity substrate as a prime-indexed orthogonal decomposition.[file:5][file:15]

\begin{lemma}[Norm decomposition]
For any \(x\in\mathcal{H}_P\) we have
\[
x = \sum_{p\in\mathbb{P}} \Pi_p x \quad\text{with}\quad \|x\|^2 = \sum_{p\in\mathbb{P}} \|\Pi_p x\|^2.
\]
\end{lemma}

\begin{proof}
Orthogonality of the ranges of \(\Pi_p\) implies that the sum is orthogonal, and by Parseval-type identities on orthogonal decompositions we have
\[
\|x\|^2 = \left\|\sum_{p\in\mathbb{P}} \Pi_p x\right\|^2 = \sum_{p\in\mathbb{P}} \|\Pi_p x\|^2,
\]
where the last equality follows from vanishing cross-terms \(\langle \Pi_p x, \Pi_q x\rangle = 0\) for \(p\ne q\).[file:15]
\end{proof}

\subsection{Prime-gated operators and block-diagonal norms}

An operator \(T\) is said to be \emph{prime-gated} if it commutes with all \(\Pi_p\):
\[
T\Pi_p = \Pi_p T \quad\text{for all } p\in\mathbb{P}.
\]

\begin{lemma}[Block-diagonal structure]
If \(T\) is prime-gated, then for each \(p\),
\[
T\Pi_p = \Pi_p T \Pi_p,
\]
and \(T\) decomposes as
\[
T = \sum_{p\in\mathbb{P}} T_p,\quad T_p := T\Pi_p = \Pi_p T \Pi_p.
\]
\end{lemma}

\begin{proof}
From \(T\Pi_p = \Pi_p T\) and idempotence of \(\Pi_p\),
\[
T\Pi_p = \Pi_p T \Pi_p,
\]
so \(T_p := T\Pi_p = \Pi_p T \Pi_p\) acts within \(\Pi_p \mathcal{H}\).[file:5][file:15]  
Sum over all primes:
\[
\sum_{p\in\mathbb{P}} T_p = \sum_{p\in\mathbb{P}} T\Pi_p = T \left( \sum_{p\in\mathbb{P}} \Pi_p \right) = TI_{\mathcal{H}_P} = T
\]
on \(\mathcal{H}_P\).
\end{proof}

\begin{proposition}[Operator norm via prime blocks]
If \(T\) is prime-gated and bounded on \(\mathcal{H}_P\), then
\[
\|T\|_{\mathrm{op}} = \sup_{p\in\mathbb{P}} \|T_p\|_{\mathrm{op}}.
\]
\end{proposition}

\begin{proof}
For the upper bound, note that for any \(x\in\mathcal{H}_P\),
\[
Tx = \sum_{p\in\mathbb{P}} T_p x = \sum_{p\in\mathbb{P}} T_p (\Pi_p x).
\]
Because the \(\Pi_p\) ranges are orthogonal and the \(T_p\) preserve those ranges,
\[
\|Tx\|^2 = \sum_{p\in\mathbb{P}} \|T_p(\Pi_p x)\|^2 \le \left(\sup_{p\in\mathbb{P}} \|T_p\|_{\mathrm{op}}^2 \right) \sum_{p\in\mathbb{P}} \|\Pi_p x\|^2
= \left(\sup_p \|T_p\|_{\mathrm{op}}^2\right) \|x\|^2.
\]
Hence \(\|T\|_{\mathrm{op}} \le \sup_p \|T_p\|_{\mathrm{op}}\).[file:5]

For the lower bound, fix any \(p_0\) and a unit vector \(x\) in \(\Pi_{p_0} \mathcal{H}\) with \(\|T_{p_0} x\| \ge \|T_{p_0}\|_{\mathrm{op}} - \varepsilon\).  
Then \(x\in\mathcal{H}_P\) and \(Tx = T_{p_0}x\), so
\[
\|T\|_{\mathrm{op}} \ge \|Tx\| = \|T_{p_0}x\| \ge \|T_{p_0}\|_{\mathrm{op}} - \varepsilon.
\]
Since \(\varepsilon>0\) is arbitrary and \(p_0\) was arbitrary,
\[
\|T\|_{\mathrm{op}} \ge \sup_{p\in\mathbb{P}} \|T_p\|_{\mathrm{op}}.
\]
Combining both inequalities gives equality.
\end{proof}

This shows that for prime-gated operators, operator norm and hence ACE budgets can be computed or bounded by analyzing each prime block separately, matching the intended ``prime-indexed recursion'' semantics.[file:5][file:15]

\section{Validation-First Gate as a Contraction Constraint}

\subsection{State counts and ACE-ledger constraint}

Let \(\mathcal{F}\) index conjecture families (e.g., prime spectra, Moonshine, EEG, CEQG).  
For each family \(F\in\mathcal{F}\), define counts:

\[
N_0^{(F)} = \#\{\text{cells in state `conjecture'}\},\quad
N_1^{(F)} = \#\{\text{cells in state `protocol-ready'}\},
\]
\[
N_2^{(F)} = \#\{\text{cells in state `running'}\},\quad
N_3^{(F)} = \#\{\text{cells in state `closed-loop'}\}.
\]

Let \(\alpha_F>0\) be a family-specific ratio and define a constitutional admissibility region
\[
\mathcal{R}_F := \left\{ (N_0^{(F)},N_1^{(F)},N_3^{(F)}) : N_0^{(F)} + N_1^{(F)} \le \alpha_F N_3^{(F)} \right\}.
\]

\begin{definition}[Validation-First admissibility]
A global state is constitutionally admissible if for all families \(F\in\mathcal{F}\),
\[
(N_0^{(F)},N_1^{(F)},N_3^{(F)}) \in \mathcal{R}_F.
\]
\end{definition}

\subsection{Projector enforcing admissibility}

Let \(\mathcal{S}\) denote the discrete state space of all cell configurations across families, and let \(\mathcal{S}_{\mathrm{adm}}\subseteq \mathcal{S}\) be the set of admissible configurations as above.[file:5]  
We can view admissibility as a projector \(\Pi_{\mathrm{adm}}\) acting on distributions over \(\mathcal{S}\).

Consider the Banach space \(l^2(\mathcal{S})\) of square-summable amplitudes over configurations, with \(\Pi_{\mathrm{adm}}\) the orthogonal projector onto the subspace supported on \(\mathcal{S}_{\mathrm{adm}}\).[file:15]

\begin{lemma}[Norm of admissibility projector]
\(\Pi_{\mathrm{adm}}\) is an orthogonal projector on \(l^2(\mathcal{S})\) and satisfies
\[
\|\Pi_{\mathrm{adm}}\|_{\mathrm{op}} = 1.
\]
\end{lemma}

\begin{proof}
By construction, \(\Pi_{\mathrm{adm}}\) acts as identity on basis elements indexed by admissible configurations and as zero on inadmissible ones; hence \(\Pi_{\mathrm{adm}}^2=\Pi_{\mathrm{adm}}\) and \(\Pi_{\mathrm{adm}}^\ast=\Pi_{\mathrm{adm}}\).[file:5][file:15]  
Orthogonal projectors on a Hilbert space have operator norm 1 by the same argument as in the proof for \(\Pi_{\mathrm{CSL}},P_E\) above.
\end{proof}

\subsection{Contraction with admissibility projector}

Let \(\widehat{C}\) denote a lifted contraction acting on \(l^2(\mathcal{S})\), representing the stochastic or deterministic update of cell configurations under the multiplicity recursion.[file:5]  
Assume \(\|\widehat{C}\|_{\mathrm{op}} \le \Lambda_m\) for some \(\Lambda_m<1\).

Define the constitutionally enforced update as
\[
\widehat{U} := \widehat{C} \Pi_{\mathrm{adm}}.
\]

\begin{proposition}[Validation-First Gate preserves contraction]
\(\widehat{U}\) is a \(\Lambda_m\)-contraction on \(l^2(\mathcal{S})\), i.e.,
\[
\|\widehat{U}f - \widehat{U}g\| \le \Lambda_m \|f-g\| \quad \text{for all } f,g\in l^2(\mathcal{S}).
\]
\end{proposition}

\begin{proof}
For any \(f,g\),
\[
\widehat{U}f - \widehat{U}g = \widehat{C}(\Pi_{\mathrm{adm}}f - \Pi_{\mathrm{adm}}g).
\]
Thus
\[
\|\widehat{U}f - \widehat{U}g\| \le \|\widehat{C}\|_{\mathrm{op}} \,\|\Pi_{\mathrm{adm}}(f-g)\|
\le \Lambda_m \|\Pi_{\mathrm{adm}}\|_{\mathrm{op}} \,\|f-g\|
= \Lambda_m \|f-g\|,
\]
since \(\|\Pi_{\mathrm{adm}}\|_{\mathrm{op}}=1\).[file:5][file:15]
\end{proof}

This formally shows that the Validation-First Gate ADR, modeled as a projector onto admissible cell-count configurations, is compatible with and preserves the contraction structure underlying ACE and lawfulness.

\section{Summary of Key Norm Relations}

For quick reference, the main operator norm bounds used in this work are:

\begin{itemize}
  \item Baseline contraction:
  \[
  \|Cx - Cy\| \le \Lambda_m \|x-y\|,\quad 0<\Lambda_m<1.
  \]
  \item ACE bound:
  \[
  \operatorname{ACE}(T) = \|(T-I)Cx^\ast\| \le \|(T-I)\|_{\mathrm{op}} \,\|Cx^\ast\|.
  \]
  \item Projector bounds:
  \[
  \|P\|_{\mathrm{op}}=1 \quad\text{for any orthogonal projector }P.
  \]
  \item Composite contraction:
  \[
  \|C\Pi_{\mathrm{CSL}} P_E\|_{\mathrm{Lip}} \le \Lambda_m.
  \]
  \item Prime-gated operator norm:
  \[
  T \text{ prime-gated} \;\Rightarrow\; \|T\|_{\mathrm{op}} = \sup_{p\in\mathbb{P}} \|T_p\|_{\mathrm{op}}.
  \]
  \item Admissibility projector:
  \[
  \|\widehat{C}\Pi_{\mathrm{adm}}\|_{\mathrm{Lip}} \le \Lambda_m.
  \]
\end{itemize}

These relations collectively justify treating ACE as a rigorous, norm-controlled observable and treating governance projectors (constitutional, ethical, admissibility) as compatible with the core multiplicity contraction dynamics.[file:5][file:15]


\nocite{*}
\bibliographystyle{unsrt}  
\bibliography{references}

\end{document}